% ============================================================
%  Georg Brandes, Samlede Skrifter. Trettende Bind (Supplementbind).
%  Kjøbenhavn: Gyldendalske Boghandels Forlag (F. Hegel & Søn), 1903.
%  Review of H. Brøchner, Bidrag til Opfattelsen af Filosofiens historiske Udvikling (1869).
%  TRANSLATION (English), printed pp. 115--123.
%  Translated from transcription.tex (Danish) in this directory, whose
%  header records how its text was established.  Method:
%  ../../../../TRANSLATION-PLAYBOOK.md.  The volume is Brandes's 1903
%  revision of pieces first printed 1865--69; the translation follows the
%  1903 text.
%
%  ---- REGISTER ----
%  Readable modern English.  Brandes's long periods are broken where English
%  will not carry them, but nothing is softened: the polemic is sarcastic,
%  and where the Danish is ironic the English is ironic.
%
%  ---- TERMINOLOGY ----
%  Aligned with texts/brandes/dualismen/translation.tex and the Nielsen and
%  Brøchner translations in this repo, since the pieces belong to the same
%  controversy:
%    Tro / Viden          = faith / knowledge   (Tro og Viden = faith and
%                           knowledge)
%    Erkendelse           = cognition (NOT "knowledge", reserved for Viden)
%    Videnskab            = science (Wissenschaft; never "natural science",
%                           for which Brandes writes Naturvidenskab)
%    Fagvidenskaberne     = the special sciences
%    absolut Uensartethed = absolute heterogeneity
%    Aand / aandelig      = spirit / spiritual;  Fornuft = reason;
%    Forstand             = understanding;  Begreb = concept
%    Samvittighed         = conscience;  Selvansvarlighed = self-responsibility
%    Livsanskuelse        = view of life;  Verdensanskuelse = world view
%    Dannelse             = cultivation;  Folkeaand = national spirit
%    Aabenbaring          = revelation;  Under = miracle
%  Danish work-titles stay in Danish, as Brandes cites them (Grundideernes
%  Logik, Om den gode Vilje, Fædrelandet, Dagbladet, ...).
%
%  ---- CONVENTIONS ----
%  - Quotation marks: the volume's single guillemets ‹...› -> ``...''.
%  - Emphasis: italic is the volume's only emphasis device; every \emph{}
%    span of the transcription is carried over as \emph{}, including the
%    Latin and other foreign phrases, which stay untranslated.
%  - Footnotes \fnmark{*)}{...} at the same anchor.
%  - The ɔ: id-est mark -> "i.e."
%  - Page markers \opage{N} and "% --- p. N ---" at the same joints as in
%    transcription.tex, so the two can be read side by side.
%  - Heads: a book under review keeps its Danish title, with an English
%    gloss in brackets; Brandes's own titles are translated.  The table of
%    contents gives the English of the volume's Indhold titles.
%
%  ---- HOW THE TRANSLATION WAS CHECKED (2026-09-22) ----
%  (1) Mechanical audit (tr_audit.py): per printed page, the same \opage
%  joints, \emph spans, heads, rules and paragraph breaks as
%  transcription.tex, and an English/Danish word ratio within 0.95--1.60:
%  9 pages, 0 flagged.  (2) Independent review: a reader who had not
%  translated these pages read every sentence of the Danish against the
%  English.  No omission found; 5 corrections applied (the genitive `dens
%  Ophavs' p. 118; `aandeligt Slægtskab' "intellectual kinship" p. 121;
%  `bevidst-ubevidst' p. 121; an added "at all" and `paa Forhaand' as "a
%  priori" p. 122).
%
%  ---- DEFECTS: DANISH AS PRINTED, ENGLISH CORRECT ----
%  The transcription is diplomatic and carries every printer's error; the
%  translation renders the SENSE and logs the divergence in a % comment at
%  the site.  The English quotation marks balance throughout.
% ============================================================

\documentclass[12pt,a4paper,openany]{book}
\usepackage[utf8]{inputenc}
\usepackage[T1]{fontenc}
\usepackage{libertinus}
\usepackage{libertinust1math}
\usepackage{textalpha}
\usepackage{graphicx}
\usepackage[english]{babel}
\usepackage[protrusion=true,expansion=true]{microtype}
\usepackage{setspace}
\usepackage[margin=1.2in]{geometry}
\usepackage{enumitem}
\usepackage{fancyhdr}
\setlength{\headheight}{28pt}

% Footnotes as the 1903 printing sets them: *), set at each site with
% \fnmark{*)}{...}, exactly as in transcription.tex.
\renewcommand{\thefootnote}{*)}
\newcommand{\fnmark}[2]{\renewcommand{\thefootnote}{#1}\footnote{#2}%
  \renewcommand{\thefootnote}{*)}}

% A piece's head, mirroring transcription.tex.  #1 = the English of the
% title as the volume's Indhold gives it (table of contents, running head);
% #2 = the head as printed, translated (a book under review keeps its Danish
% title, glossed in brackets).  The short rule under every head is printed.
\newcommand{\piecehead}[2]{\clearpage\phantomsection
  \addcontentsline{toc}{section}{#1}\markboth{#1}{#1}%
  \begin{center}\large #2\\[6pt]\rule{2cm}{0.4pt}\end{center}\par\medskip}
\newcommand{\piecerule}{\par\begin{center}\rule{1.5cm}{0.4pt}\end{center}\par}

\usepackage{hyperref}
\hypersetup{pdftitle={Brandes, Collected Writings XIII: The Historical Development of Philosophy (1869)}, pdfauthor={Georg Brandes}, hidelinks}

\pagestyle{fancy}
\fancyhf{}
\fancyhead[LE,RO]{\thepage}
\fancyhead[RE]{\textit{Collected Writings XIII (1903)}}
\fancyhead[LO]{\textit{\leftmark}}
\setstretch{1.3}

\title{\textbf{The Historical Development of Philosophy (1869)}\\[8pt]\large Georg Brandes, \textit{Samlede Skrifter} [Collected Writings], vol.~XIII (Supplementbind),\\ pp.~115--123}
\author{}
\date{Copenhagen: Gyldendal (F.\ Hegel \& Søn), 1903\\[10pt]
  \small Translated from the Danish}

% ---- original-page markers ----
% \opage{N} marks where printed page N begins, at the same joint as in
% transcription.tex (mid-sentence where the page turns mid-sentence).
\usepackage{marginnote}
\newcommand{\opage}[1]{\marginnote{\footnotesize\textit{[#1]}}}
\newcommand{\apage}[1]{\marginnote{\footnotesize\textit{[#1]}}}
\begin{document}
\maketitle
\thispagestyle{empty}
\tableofcontents

\mainmatter

% --- p. 115 ---
\piecehead{The Historical Development of Philosophy}{\opage{115}\emph{H. Brøchner: Bidrag til Opfattelsen af Filosofiens}\\
\emph{historiske Udvikling}\\[2pt]{\normalsize[Contributions to the Understanding of the Historical Development of Philosophy]}\\[2pt](1869)}
One generally reads with a certain astonishment that the most
significant men of the past have most often been precisely those
to whom their contemporaries paid the least heed; one is surprised
to learn which men were preferred to them, what current of the
age they had against them, how solitary their position was; and
one feels assured that in our day such things are over, that our
age knows how to find merit in its most hidden corner, seize it
and draw it forth into the light; one flatters oneself with the
thought that the names which sound most often and most loudly in
our day will also be the ones on the lips of posterity. Not many
things in the world are as certain as that in this one is sadly
mistaken. Names that are now celebrated, persons whom everybody
now knows, will be a nothing to posterity, and personalities who
are now known and valued by a small circle will leave behind them
lasting traces of their life and work. But worse off, after all,
is he who outlives his fame than he whose fame grows so slowly
that only when the age has caught up with the personality does it
stand out in clear outline.

Professor \emph{Brøchner} is a man who belongs to the latter
class; he is no popular man; he has never taken a single step
toward popularity, and he is far too far ahead of his surroundings
to be valued by them; for one values
% --- p. 116 ---
\opage{116}only what one understands. He has worked in silence, and had he
been carried off by illness before he at last, last year, began
the publication of his writings, scarcely half a score of people
would have known that the country had lost its most thorough and
most incisive thinker. Fortunately he has come through his long
illness, and the publication of his works, awaited with such great
longing by his few disciples, has been begun.

When one considers the book whose title is given above, one's
first impression is a painful feeling. Is it not sad that Denmark
has gradually been reduced to such smallness that the best
writings can find no readers, can reap no recognition, can yield
no return and can receive no criticism? Would it not be better,
both for the man of science himself and for the whole of European
cultivation, if men like Professor Brøchner and our other
significant men of science lived in one of the great principal
countries, were developed through criticism and received stimulus
from every side? Perhaps it would be better. But it seems to me
that the lot of such a man has its beauty even under the present
circumstances. He and his like are the farthest advanced vedettes
or outposts of the great European culture in the high North; their
position is the most exposed; for outpost duty is severe and
dangerous; but precisely thereby their scientific life demands a
firmness of character, a courage, at certain times a soldier's
daring, which less uncommon conditions do not develop, and which
inspires
% Danish prints `kans Talent' for `hans Talent'; rendered as the sense
% ("his talent").
love for the man, adorns and strengthens his talent.

Professor Brøchner's work as a teacher at the University,
though it has lasted several years, is not yet of such long
standing that its influence can be clearly traced; but that it is
significant is certain, and anyone who knows many of the younger
students will be able to bear witness to it. Professor Brøchner
is the first philosopher freed from all theology that our
University has possessed; if G. Sibbern, who has had to pay
dearly for his liberal-mindedness, is excepted, he is also the
only living Danish one. He was appointed at a time when
theological influences of the most varied kinds were powerful here
at home; he has spoken from his solitary lectern to his not always
numerous audience, while all around him the spiritually
reactionary voices that drowned out everything else shouted to one
another and confused every thinking brain. He has held fast to the
conviction which from his early youth showed
% --- p. 117 ---
\opage{117}itself to him in the clarity of truth; in the strongest
anti-rational storms he has held the banner of reason fast.

How many of the University's other professors, e.g. of the
Faculty of Natural Science, have not shared his way of seeing,
and yet, as is the common custom among learned people, have kept
quiet in order to avoid opposition and unpleasantness, and have
always maintained that there was no reason to speak ``so long as
no one did anything to them, made no encroachment on their
freedom.'' As if a man were a man of science if the truth did not
burn him like a live coal on the tongue; as if he who kept silent
did not leave in the lurch the individual who speaks, and did not
contribute his share to producing the false appearance that the
latter, with his way of seeing, stands literally alone, that his
view of life is a kind of whim, an unreasonable oddity, sprung up
in a brain twisted by false metaphysics, etc. As I said, Professor
Brøchner has found no support. But this much is certain: if there
is anyone among the younger men who study the sciences of the
spirit who in the course of his university career has received a
strong impression of the unconditional right of free inquiry, of
the validity of humanity over against all dogmas, then he surely
owes this impression to Professor Brøchner; and the reviewer can
at least testify for himself that scarcely anything, amid the
bewildering battle of crisscrossing opinions and systems, has been
such a support to him and such a direction to seek the truth by
the way of reason as Professor Brøchner's strictly scientific
inquiry and manfully unbending stance.

The treatise on \emph{Filosofiens historiske Udvikling} is surely
the most significant of the writings the author has so far
published. It is a work that will keep its place in our
literature. It will be an excellent guide for anyone who takes up
the study of the history of philosophy, and it is so new and
original, so well thought through, so tested again and again, that
even in the rich German literature it would attract great notice
--- what a role, then, does its publication not play in a
philosophical literature as poor as the Danish!

Professor Brøchner has not wished to give the history of
philosophy in and for itself. He has set himself a larger goal:
he has wished to give the philosophy of this history. What does
that mean? When does history get its philosophy? It gets it when
it is satisfactorily shown that it is not a sequence of
accidental, disconnected,
% --- p. 118 ---
\opage{118}lawless events, when, rather, it is seen to be governed by law. That
it must be subject to a lawfulness Brøchner establishes by
pointing to the essence of its originator, the human race, whose
unfolding cannot be accidental or planless but must be a real
development. In this word \emph{development}, the concept that is
the center of the whole of recent German and English philosophy,
there lie rolled up all the determinations or laws that a thinking
consideration of history is able to unfold. If history has a
coherence, then no historical state of affairs, of any kind that
can be named, can be without lateral influence in all directions;
if history has a development, then no historical fact, no matter
of what kind, can be without after-effect out into the future. If
one takes a single small domain out of history, literary history
for example, it stands in connection with all the other branches;
if one takes out of that in turn a single section such as the
history of philosophy, one once again artificially isolates a
single part of history, which forms a whole of its own only
relatively. But the development that takes place in the great life
of culture is found also in this particular organ of it, just as,
conversely, the lawfulness we find here is a guarantee to us that
arbitrariness cannot reign at any point.

The individuals of whom the race consists no doubt believe
that they think as their own individual reason prompts them, as
they can and as they will. In reality, however, they think also,
and first and foremost, as they \emph{must}. They are derived
beings; they have their thoughts from themselves as little as
their life; just as the universal, just as the race \emph{lives}
in them, so too the human race \emph{thinks} in the individuals.
The ideas it is preparing appear, first, in an embryonic state in
individual thinkers, but without these being conscious of it
themselves, a good while before those ideas, fully developed, take
their place in science. Second, the race's stage of development at
the given point in time stands at the same time like a fate over
the individuals. Although the very ideas whose organs the
individuals are urge them in many ways to draw further
consequences from those thoughts, to pursue them out in new,
future directions, the spirit of the age is nevertheless a power
that encompasses all the individual's endeavors, holds him back
without his being conscious of it, and prevents him from pursuing
a thought beyond the horizon of his age.

% --- p. 119 ---
\opage{119}Yet the fate that governs the individual's thinking is more
composite still: not only is he unable to overstep the bounds of
the contemporary stage of the race, but so long as this stage has
not yet fully ripened and come to rest in equilibrium about its
own center, the preceding, abandoned stage still constantly works
as a power that hampers the individual's striving. This too is a
law, grounded in the concept of development as well as confirmed
by historical experience. If, finally, one notices that often
within one and the same historical age diametrically opposed
tendencies assert themselves, then the philosopher of history is
able even here to show us how the standpoint of development on
which the whole age stands is, without the two opposed camps being
conscious of it, common to them both, and unites those who regard
themselves as opponents in principle under the all-governing form
of development, which, marking friend and foe alike with its
stamp, shapes itself into two opposed schools as into two
hemispheres that complete each other.

It is laws of this nature that Professor Brøchner has
demonstrated in the history of philosophy, and in order to confirm
the results of his inquiry still further he has, as is his habit,
followed the procedure of furnishing both proof and counterproof.
He does not content himself with illustrating his line of thought
by examples; but once this is done he chooses such cases as
apparently constitute exceptions to the law, and shows how it
finds its confirmation in them as well.

Yet the excellence of his new book is not to be sought solely
in the demonstration of laws, nor in the study of the relation
between the conscious and the unconscious in human nature, a
domain which, incidentally, may very nearly be called Professor
Brøchner's special domain; it rests besides on the exhaustive
knowledge of the history of philosophy by which the author takes
so high a rank among the thinking minds of Europe. It is this that
has made it possible for him to furnish such masterly
investigations as those of Aristotle and Kant, who are here
understood in a depth and with a clarity found scarcely anywhere
else. So clear a line of thought begets a clear and beautiful
exposition, and there are therefore large parts of this work where
the author has succeeded, despite the refractory nature of the
subject, in bending it into harmonious forms, and, despite the
refractory nature of his pen, now and then in
% --- p. 120 ---
% Danish prints a comma before the capital of a new sentence (`Farve,
% Saaledes'); rendered as the sense, with a full stop.
\opage{120}letting the mood lend the presentation a gentle, spiritualized
color. So it is, for example, in the portrayal of Neoplatonism and
of Giordano Bruno.

This is not the place to discuss the details of the book, even
if the reviewer were equal to the task, which he is not. I shall,
on the other hand, raise some objections to the plan of the book
and to the author's standpoint.

The plan is not altogether even, and the author's way of seeing
is not the same toward all subjects. He has, of course, the right
to present the longitudinal coherence of philosophy by itself and
to see it as a relative whole. But it is also clear that such a
conception must of necessity become very abstract and rather
one-sided. What seems to me somewhat unsteady is this, that the
author sometimes takes the lateral coherence into account and
sometimes does not. In my view no philosophy can be fully
understood unless it is seen in its close connection with all the
other branches of the literature of its time, with the
contemporary historical events
% Danish prints a comma before a new sentence (`Folkeaand, Naar');
% rendered as the sense, with a full stop.
and with the national spirit of its people. When the subject is
Greek philosophy, the author is willing to see it in its
connection with the whole Greek spirit and culture, and to hold
that only in this way can it be understood. But when the subject
is modern French, German, English philosophy, it is treated as
though it led a secluded existence without relation to the
national spirit, without relation to the country's poetry,
history, etc. Is it not strange not to introduce the presentation
of German philosophy at all with a sketch of the gifts and limits
of the German spirit, but to regard it without further ado as
all-embracing? It perhaps seems so to Brøchner because he stands
so near it; but is sheer narrow-mindedness, then, to blame for its
never having caught on in England? I doubt it in the very highest
degree. I for my part hold not only that at bottom it is
unintelligible without a preliminary study of the German national
character through its poetry, its criticism, its literature as a
whole, but also that several distinctive traits in its individual
heroes, as for example the strong assertion of the I, the strong
enthusiasm for freedom in Fichte, are first seen in their true
light when one shows Napoleon's despotism over Germany in the
background.

I admire Brøchner's gift for presenting the main points and the
main features of a system; but in all philosophers is
% --- p. 121 ---
\opage{121}the systematic element what is most essential and most important?
Is the skeleton everything, the body nothing? Ought one not at
least to be wary of stamping interesting, heterogeneous
philosophical phenomena with the same official school-label when
quite different national spirits have found expression in them?
% Danish prints `fremstilies' for `fremstilles'; rendered as the sense
% ("is presented").
Condillac is usually presented as an appendix to Locke; is there
not all the same a yawning gulf between their whole manner of
philosophizing? Is there not, on the score of nationality, a
still stronger intellectual kinship between Descartes and Condillac
than there is, with respect to the results of thought, between
Condillac and Locke? Does the key to Descartes's philosophy not
lie in the spirit of the whole great French century, and is it
not connected with the theology, the poetry, the legislation and
the landscape gardening of the time? In a word, do not the spirit
of the stock and the national spirit also belong to the fate that
is mightier than the fate of individuals, and that works on them
unconsciously and hampers them? It seems to me that this new fate
deserved a stronger emphasis.

And since we are speaking of fate, has not Brøchner himself his
own? He belongs to the large group of philosophers who since
Hegel's time have tried, while retaining everything the great
master won, to make good his mistakes, avoid his overreachings,
and immerse themselves in the reality of which he, with his
contempt for nature, often disdained empirical knowledge. Such
points as the inquiry into laws and the demonstration of the
relation between the conscious and the unconscious show how fully
Brøchner stands at the height of his age and shares in its
endeavors. Nevertheless, with all this he still stands, after all,
on Hegel's ground, and I will not deny that I am inclined to % `den Tro': ordinary belief here, not the controversy's `faith'
believe that a conscious-unconscious after-effect of the % `bevidst-ubevidst': Brandes turns Brøchner's own theme (the conscious and the unconscious) on him
panlogism he has abandoned sometimes holds him back and sometimes
misleads him into a priori constructions that scarcely hold water.
Among these I count his division of philosophy into three main
epochs, of which two are past, while the third, the philosophy of
the future, is now about to begin. The task that he sets this
philosophy of the future the thinkers of the preceding epoch have
already striven to accomplish, and whether the thinkers who arise
in the centuries immediately following will be more fortunate than
those of the present one in avoiding a one-sided emphasis on
nature or on spirit is a
% --- p. 122 ---
\opage{122}question. Whether they will in the end make use of the method
of inquiry that Brøchner prescribes for them is a still greater
one.

The dialectical method, which has so far been practiced only in
a single country and over a short span of time, has not stood the
test in its original form. It has been abandoned even in Germany,
and it is not used in any other country. That the dialectical is
a power in the human spirit and in the universe I am far from
doubting; but with what adaptation the method can be used Brøchner
has not shown. All of us who work would like nothing better than
an excellent tool; but since the highly ingenious one that has so
far been put into our hands cannot be used, no one can hold it
against us that for the present we prefer to make use of the coarser
implements of the special sciences, which at least do not break in
the hand.

Finally, I see a remnant of Hegel's influence in the way in
which Brøchner disposes of Oriental philosophy and defines the
Greek. The Oriental he does not reckon with --- a little
prematurely, it seems to me. There exist extremely complex, % `en Smule paa Forhaand': in advance, before examination; cf. `Forhaands-Konstruktioner', p. 121
extraordinarily interesting Indian systems which have not yet been
translated at all; they are, if the reviewer may take others at
their word, philosophical in their essence; they have no special
relation to religion, and they must surely be examined before they
can be sacrificed to the form of philosophy which, as representing
``the immediate unity of nature and spirit,'' must necessarily be
the first.

Brøchner arrives at this last determination by beginning with
stamping the Greeks as \emph{the people of beauty}. This too seems
to me an undoubted after-effect of Hegelianism. Such propositions
are all too true, that is to say, all too abstract to be true. The
Greeks are no more the people of beauty than the Jews are the
people of religion, or \emph{Don Juan} the perfect opera, etc. The
Greeks express a single, definite, historical ideal of beauty, and
they came by a purely historical route to express it. This ideal
is only in certain respects the highest that has arisen up to now,
and it became so for purely historical reasons. The Greeks are the
inventors and founders of cities that are at the same time states.
The perilous situation of these cities among only half-subjugated
tribes, and their life of feuding among themselves, made necessary
the gymnastic education, which in turn brought it about that the
perfect gymnast was valued and admired to an unheard-of degree,
% --- p. 123 ---
\opage{123}was crowned with wreaths at public festivals, won the right to a
statue, and became the ideal of beauty. Is not, then, the
determination that the Greeks are the people of the city-state
more deep-lying and less unconditional than the determination that
they are the people of beauty? It seems to me that it is, and it
seems to me reasonable that from the Greeks' wholly peculiar
political life one would be able to arrive at many of the
determinations that are the most decisive for their philosophy.

It is possible that I am wrong in some of these objections. If,
however, they are not unfounded, the defects which they presumably
hit nonetheless in no way rule out the great excellence of the
book. At bottom they only establish the validity of one of the
author's own laws, namely this: that the older epoch, so long as
the new one has not yet emerged in its full development, works,
unconsciously for the individual thinker, to a certain degree as a
hindrance on his conscious striving to bring forth the philosophy
of the new age.
\piecerule

\end{document}
